\documentclass{article}

\usepackage{geometry}
\usepackage{makecell}
\usepackage{array}
\usepackage{multicol}
\usepackage{setspace}
\usepackage{changepage}
\usepackage{booktabs}
\usepackage[explicit]{titlesec}
\usepackage{hyperref}
\usepackage{graphicx}
\usepackage{cprotect}
\usepackage{float}
\newcolumntype{?}{!{\vrule width 1pt}}
\newcommand{\paragraphlb}[1]{\paragraph{#1}\mbox{}\\}
\newcommand{\subparagraphlb}[1]{\subparagraph{#1}\mbox{}\\}
\renewcommand{\contentsname}{Inhaltsverzeichnis:}
\renewcommand\theadalign{tl}
\setstretch{1.10}
\setlength{\parindent}{0pt}

\titleformat{\section}
  {\normalfont\Large\bfseries}{\thesection}{1em}{\hyperlink{sec-\thesection}{#1}
\addtocontents{toc}{\protect\hypertarget{sec-\thesection}{}}}
\titleformat{name=\section,numberless}
  {\normalfont\Large\bfseries}{}{0pt}{#1}

\titleformat{\subsection}
  {\normalfont\large\bfseries}{\thesubsection}{1em}{\hyperlink{subsec-\thesubsection}{#1}
\addtocontents{toc}{\protect\hypertarget{subsec-\thesubsection}{}}}
\titleformat{name=\subsection,numberless}
  {\normalfont\large\bfseries}{\thesubsection}{0pt}{#1}

\hypersetup{colorlinks,
    		citecolor=black,
    		filecolor=black,
    		linkcolor=black,
    		urlcolor=black}

\geometry{top=12mm, left=1cm, right=2cm}
\title{\vspace{-1cm}Projektmanagement}
\author{Andreas Hofer}

\begin{document}
	\maketitle
	\tableofcontents
	\section{System Design}
	Systems Engineering ist eine Methode zur Projektplanung, welches heutzutage in einem Großteil der Welt angewandt wird. Dabei wird angenommen, dass das Problem ein System ist, welches in dessen Elemente gespalten und in Beziehung zueinander gebracht werden soll. Unter einem Problem versteht man allgemein die Differenz zwischen einem IST- und einem SOLLzustand, also einer Situation wie sie in diesem Moment besteht und man eine Idee hat wie der Idealzustand aussehen soll.
	\subsection{Routine- und Pionierprobleme}
	Probleme werden weiterhin in Routine- und Pionierprobleme geteilt. Routineprobleme sind bereits gelöst oder haben einen gut definierten Lösungsweg um zur Lösung zu finden. Pionierprobleme hingegen haben keinen genauen Lösungsweg oder er ist nur ungenau definiert. Da Pionierprobleme schwerer zu lösen sind, wird auch ein Großteil des Fokus von Systems Engineering darauf gesetzt.
	\begin{itemize}
		\item{Methodik}
		\item{Psychologie und emotionale Intelligenz}
		\item{Ethik und Moral}
		\item{Fachwisen und Erfahrung}
		\item{Situationskenntnisse}
	\end{itemize}
	Während jeder dieser Säulen wichtig ist, besteht in dieser Lehrveranstalung ein Fokus auf die Methodik. Diese ist die Wegleitung um den Lösungsweg definieren zu können. Systems Engineering ist ein Bauplan um die Methodik eines Problems definieren zu können.
	\subsection{Systems Engineering}
	Ein System hat stets definierte Komponenten:
	\begin{itemize}
		\item{Elemente}
		\begin{itemize}
			\item{Bausteine eines Systems}
		\end{itemize}
		\item{Beziehungen}
		\begin{itemize}
			\item{Die Abhängigkeiten der Elemente innerhalb des Systems}
		\end{itemize}
		\item{Umwelteinflüsse}
		\begin{itemize}
			\item{Einflüsse auf das System welche selbst nicht Teil des Systems sind}
		\end{itemize}
		\item{Systemgrenzen}
		\begin{itemize}
			\item{Was das System von der Umwelt trennt}
		\end{itemize}
		\item{Systemziel}
		\begin{itemize}
			\item{Das was man mit dem System erreichen will}
		\end{itemize}
	\end{itemize}
	\subsubsection{Vorgehensmodell}
	Beim Vorgehensmodell wird definiert wie man überhaupt sein System definieren will. Dabei sind vier Grundgedanken relevant:
	\paragraphlb{Vom Groben zum Detail}
	Da man nicht sofort eine detaillierte Angabe des Problems und dessen Ziels angeben kann, sollte man zuerst eine grobe Vorstellung der Situation und des Ziels definieren. Anhand dieser Basis definiert man seine Vorstellungen dann feiner um so iterativ zum Endergebnis zu kommen.
	\paragraphlb{Denken in Varianten}
	Eventuell gibt es nicht sofort einen offensichtlichen Lösungsweg, sondern mehrere potenzielle Varianten. In der Regel sollte man drei verschiedene Varianten von Lösungsansätzen definieren.
	\paragraphlb{Projektphasen}
	Bei Projektphasen unterteilt man die Methodik in verschiedene Phasen, welche danach in Serie abgearbeitet werden kann:
	\begin{enumerate}
		\item{Vorstudie}
		\begin{itemize}
			\item{Es werden im Vorhinein mögliche Lösungswege definiert (idealerweise 3)}
			\item{Es soll im Vorhinein das Problem verstanden werden}
		\end{itemize}
		\item{Hauptstudie}
		\begin{itemize}
			\item{Anhand der möglichen Lösungswege werden konkrete Aktionen (idealerweise wiederum 3) erkundet und der beste Weg gewählt}
		\end{itemize}
		\item{Detailstudie}
		\begin{itemize}
			\item{Der gewählte Weg wird detaillierter für verschiedene Wege überlegt}
		\end{itemize}
		\item{Implementation}
		\begin{itemize}
			\item{Der Detailplan wird in die Tat umgesetzt}
		\end{itemize}
	\end{enumerate}
	\paragraphlb{Problemlösungszyklus}
	Der Problemlösungszyklus (PLZ) beschreibt die Gliederung der Methodik in einzelne Projektphasen:
	\begin{itemize}
		\item{Zielsuche}
		\begin{itemize}
			\item{"Wo stehen wir?"}
			\item{"Wo wollen wir hin?"}
		\end{itemize}
		\item{Lösungssuche}
		\begin{itemize}
			\item{"Welche Lösungen gibt es?"}
			\item{"Welche Lösungen kommen in Frage?"}
		\end{itemize}
		\item{Auswahl}
		\begin{itemize}
			\item{"Welche Lösung ist die Beste?"}
			\item{"Welche Lösung können wir empfehlen?"}
			\item{"Für welche Lösung entscheidet sich der Auftraggeber?"}
		\end{itemize}
	\end{itemize}
	\subsection{Problemlösungsprozess}
	\subsubsection{Situationsanalyse}
	Der erste Schritt ist die Analyse der gegebenen Situation. Bei dieser Analyse durchläuft man vier Phasen:
	\begin{enumerate}
		\item{Problemfeldabgrenzung}
		\item{Umwelt- und Systemanalyse}
		\item{Ursachen ergründen}
		\item{Ergebnis darstellen}
	\end{enumerate}
	\paragraphlb{Problemfeldabgrenzung}
	Bei der Problemfeldabgrenzung definiert man welche Elemente im System eine Relevanz besitzen. Dabei kann man grob in drei Schritten vorgehen:
	\begin{enumerate}
		\item{Problemfeld definieren}
		\begin{itemize}
			\item{Alle Elemente welche für das System relevant sind, egal ob System- oder Umweltelemente, werden definiert.}
		\end{itemize}
		\item{Beziehungen definieren}
		\begin{itemize}
			\item{Die wichtigen Beziehungen der einzelnen Elemente zueinander aufbauen}
		\end{itemize}
		\item{Elemente kategorisieren}
		\begin{itemize}
			\item{Elemente in System- und Umweltelemente teilen. Die vorherrschende Frage ist dabei: "Habe ich Einfluss auf das Element oder beeinflusst das Element mich?"}
		\end{itemize}
		\item{Syst}
	\end{enumerate}
	\paragraphlb{Umwelt- und Systemanalyse}
	Jetzt hat man zwar definiert, welche Elemente das System beeinflussen, muss jedoch nun herausfinden, was diese Faktoren wirklich sind. Wenn man zuvor zum Beispiel ein Budget als Element festgelegt hat, muss man nun bestimmen wie hoch dieses Budget eigentlich ist. Diese Informationen können auf mehrere Wege beschafft werden, unter anderem:
	\begin{itemize}
		\item{Literatur- und Internetrecherche}
		\item{Interviews}
		\item{Fragebögen}
		\item{Beobachtungen}
	\end{itemize}
	Diese gesammelten Informationen müssen danach verarbeitet und aufbereitet werden. Eine Strategie hierfür ist die \textit{SWOT-Analyse}
	\paragraphlb{SWOT-Analyse}
	SWOT steht für:
	\begin{itemize}
		\item{Stärke (Strength)}
		\item{Schwäche (Weakness)}
		\item{Chance (Opportunity)}
		\item{Gefahr (Threat)}
	\end{itemize}
	Womit die Gegebenheiten für jedes Element analysiert werden können.
	\paragraphlb{Zielformulierung}
	Die SWOT-Analyse hilft dabei Ziele für das Projekt zu definieren. Diese Ziele sollten vier Dimensionen erfüllen:
	\begin{enumerate}
		\item{Zielinhalt}
		\begin{itemize}
			\item{Was soll passieren?}
		\end{itemize}
		\item{Zielausmaß}
		\begin{itemize}
			\item{In welchem Ausmaß soll etwas erreicht werden?}
		\end{itemize}
		\item{Zielregion}
		\begin{itemize}
			\item{Wo soll das Ziel erreicht werden?}
		\end{itemize}
		\item{Zielfristigkeit}
		\begin{itemize}
			\item{Bis wann soll das Ziel erreicht werden?}
		\end{itemize}
	\end{enumerate}
	Doch nicht jedes Ziel ist gleich wichtig. Man unterscheidet stets zwischen Musszielen, Ziele die erreicht werden müssen, und Wunschzielen, Ziele die nützlich sind, wenn sie erreicht werden, jedoch nicht kritisch. \\
	Man muss ebenfalls darüber nachdenken, wie diese Ziele miteinander interagieren. Ziele können einander komplementieren, also hilft die Erfüllung eines Ziels dabei ein anderes Ziel zu erreichen. Ziele können jedoch auch in Konflikt stehen, wodurch die Erfüllung eines Ziels die Erfüllung eines anderes Ziels erschweren. Ziele können jedoch auch neutral zueinander stehen, wenn sie nicht miteinander interagieren.
	\paragraphlb{Synthese/Analyse}
	Bei der Synthese und Analyse sollen nun anhand der Situationsanalyse und dessen Zielführung konkrete Lösungwege erarbeitet werden um das Ziel zu erreichen. Dabei stellt die Synthese ein Brainstorming dar, bei welchem mögliche Ideen ohne Überlegung ihrer Machbarkeit niedergeschrieben werden. \\
	Erst bei der Analyse beobachtet man jeden Lösungsweg kritisch um zu eruieren, ob dieser realistisch ist. Im Endeffekt sollten drei konkrete und realistische Lösungswege vorgelegt werden.
	\paragraphlb{Bewertung}
	Bei der Bewertung sollen diese drei potentiellen Varianten nun verglichen werden um danach die Beste zu finden. Bei der Bewertung unterscheidet man zwischen qualitativen und quantitativen Bewertungsverfahren.
	\subparagraphlb{Qualitativ}
	Qualitative Bewertungsverfahren können nicht in Zahlen abgebildet werden. Bei qualitativen Verfahren zählt eine verbale Einschätzung oder eine Liste an Checklisten. Solche Bewertungen sind oft subjektive Einschätzungen oder binäre Kriterien, welche entweder existieren oder nicht.
	\begin{enumerate}
		\item{Verbale Enschätzung}
		\begin{itemize}
			\item{Bei der verbalen Einschätzung werden Pro/Contra Argumente gegenübergestellt und anhand dieser wird danach eine Entscheidung getroffen.}
			\item{Sehr subjektiv, da die Wichtigkeit der einzelnen Kriterien nicht ersichtlich ist. Aus diesem Grund sollte eine verbale Einschätzung in der Praxis nicht verwendet werden.}
		\end{itemize}
		\item{Checkliste}
		\begin{itemize}
			\item{Die Checkliste ist eine Weiterentwicklung der verbalen Einschätzung. Hierbei werden Kriterien definiert und jede der Varianten danach anhand dieser Bewertet.}
			\item{Die Checkliste lässt jedoch auch keine Abstufung der Wichtigkeit zu und benötigt oft schwammige Aussagen wie "Wenig", "Mittel" oder "Viel"}
		\end{itemize}
		Um diese Probleme zu vermeiden sollten Bewertungskriterien stets gewichtet werden. So kann man stets feststellen wie wichtig ein Kriterium im Vergleich der anderen ist. Ein Weg um diese Gewichtungen zu ermitteln ist ein Paarvergleich. Dabei werden jeweils zwei Kriterien gegenübergestellt, und eine als wichtiger und eine als unwichtiger gestuft.
	\end{enumerate}
	\subparagraphlb{Quantitativ}
	Quantitative Bewertungsverfahren lassen sich in Zahlen abbilden und statistisch auswerten. So kann man mit arithmetischen Verfahren die beste Entscheidung abbilden.
	\subparagraphlb{Kombination der Verfahren}
	Diese beiden diskreten Verfahren kann man auch kombinieren um die Vorteile beider verwenden zu können. Zwei Verfahren zur Kombination sind die Nutzwertanalyse und Kosten-Wirksamkeits-Analyse. \\
	{\bold \Large Nutzwertanalyse} \\
	Die Nutzwertanalyse versucht den menschlichen Entscheidungsprozess abzubilden. Dabei werden Bewertungskriterien aufgelistet und die Einzelnen Varianten aufgelistet. \\
	Danach wird für jede der Kriterien eine unabhängige Notenskala aufgestellt. Dabei orientiert man sich nicht an den verfügbaren Werten der Varianten sondern an den eigenen Erwartungen. Dabei erhält ein unakzeptabler Wert die Note 0 und der Idealwert die Note 10. Danach werden diese unabhängig definierten Werte den Werten der Varianten gegenübergestellt. Der einfachste Weg um den Werten eine Note zuzweisen, ist eine lineare Interpolation: Der gesuchte Wert $x_1/x_2/x_3$ geteilt durch die Basis der Skala ist gleich dem oberen Ende der Notenskala geteilt durch die Differenz der Wertskala: $\frac{x_1}{Basis}=\frac{M_{Note}}{M_{Werte}}$. Wenn die Werte also von 40 bis 130 gehen, und die Notenskala von 0 bis 10 geht, ist die Basis stets der Wert minus dem unteren Ende der Werte: $y_1-40$ und der Maximalwert 10 und $130-40$: $\frac{10}{90}$. Um also $x_1$ zu finden muss man vom Wert der Variante $y_1$ 40 abziehen und dem Quotienten von Note und Wert gegenüberstellen: $\frac{x_1}{y_1-40}=\frac{10}{90}$ \\
	{\bold \Large Kosten-Wirksamkeits-Analyse}
	Ein großese Problem der Nutzwertanalyse ist, dass jede der Kriterien gleich behandelt wird. Aus diesem Grund versucht die Kosten-Wirksamkeits-Analyse Kriterien in Wirksamkeits- und Kostenkriterien zu unterteilen. \\
	Der Prozess um die Kosten-Wirksamkeits-Analyse durchzuführen ist zuerst die Nutzwertanalyse wie normal durchzuführen, dabei jedoch alle Kriterien mit Kosten zu entfernen (Dabei nicht vergessen die Kriterien neu zu gewichten). Danach werden alle Kosten aufsummiert und durch die Nutzwerte der Nicht-Kostenkriterien geteilt.
	\section{Projektmanagement}























  
\end{document}